\documentclass{amsart}

\usepackage{amsfonts,amssymb,amsmath,enumerate,verbatim,mathtools,tikz,bm,mathrsfs,tikz-cd,hyperref,comment,stmaryrd,amsthm}
\usepackage{colonequals}
\usepackage{adjustbox}
\usepackage[all,pdf]{xy}
\hypersetup{
    colorlinks=true,
    linkcolor=blue,
    filecolor=magenta,      
    urlcolor=cyan,
}
\usepackage{cleveref}
\usepackage{adjustbox}
\usepackage{leftindex}

\makeatletter
\@namedef{subjclassname@2020}{\textup{2020} Mathematics Subject Classification}
\makeatother





\author[Jian Liu]{Jian Liu}
\address{School of Mathematics and Statistics, and Hubei Key Laboratory of Mathematical Sciences,  Central China Normal University,  Wuhan 430079, P.R. China}
\email{jianliu@ccnu.edu.cn}
%\urladdr{\url{https://orcid.org/0000-0001-8360-7024}}


\keywords{dominant local ring, descent, singularity category, homotopy category, acyclic complex, injective module, finite CM type}
\subjclass[2020]{18G80 (primary); 13C11, 13D09 (secondary)}


\newcommand{\com}[1]{{\color{green} #1}}
\newcommand{\old}[1]{{\color{red} #1}}
\newcommand{\new}[1]{{\color{blue} #1}}


\newcommand{\sfS}{\mathsf{S}}
\renewcommand{\S}{{\mathcal{S}}}
\DeclareMathOperator{\cone}{cone}
\DeclareMathOperator{\rank}{rank}
\DeclareMathOperator{\spn}{span}
\DeclareMathOperator{\amp}{amp}
\DeclareMathOperator{\embdim}{embdim}
\DeclareMathOperator{\dg}{dg}
\DeclareMathOperator{\codepth}{codepth}\DeclareMathOperator{\depth}{depth}
\DeclareMathOperator{\chr}{char}
\DeclareMathOperator{\hh}{HH}
\DeclareMathOperator{\h}{H}
\DeclareMathOperator{\bb}{B}
\newcommand{\n}{\mathfrak{n}}
\newcommand{\lb}{\llbracket}
\newcommand{\VS}{\mathsf{VS}}
\newcommand{\rb}{\rrbracket}
\newcommand{\athir}[2]{\displaystyle \bigoplus_{\sigma^{#1}\in C_{#1}}K_{#2}    }
\newcommand{\surj}{\twoheadrightarrow}
\newcommand{\CC}{\mathcal{C}}
\newcommand{\RR}{\mathcal{R}}
\newcommand{\VV}{\mathcal{V}}
\newcommand{\PV}{\mathbb{V}}

\newcommand{\F}{\bm{F}}
\newcommand{\FF}{\mathcal{F}}
\newcommand{\EE}{\mathcal{E}}
\newcommand{\Q}{\mathbb{Q}}
\newcommand{\Z}{\mathbb{Z}}
\newcommand{\PS}{\mathbb{P}}
\newcommand{\B}{\mathcal{B}}
\newcommand{\R}{\mathbb{R}}
\newcommand{\A}{\mathcal{A}}
\newcommand{\AI}{\mathbb{A}}
\newcommand{\sfA}{\mathsf{A}}
\newcommand{\T}{\mathsf{T}}
\newcommand{\G}{\mathsf{G}}
\newcommand{\gt}{{\rm gentime}}
\newcommand{\Rdim}{{\rm Rdim}}


\newcommand{\D}{\mathsf{D}}
\newcommand{\K}{\mathsf{K}}
\newcommand{\C}{\mathsf{C}}
\newcommand{\N}{\mathbb{N}}
\newcommand{\ov}[1]{\overline{#1}}
\newcommand{\vp}{\varphi}

\newcommand{\kos}[2]{{#1}/\!\!/{#2}}
\newcommand{\para}{\mathbin{\!/\mkern-5mu/\!}}

\newcommand{\con}{\subseteq}
\newcommand{\Aff}{\mathbb{A}_n^k}
\pgfdeclarelayer{bg}    % declare background layer
\pgfsetlayers{bg,main} 
\newcommand{\x}{{\bm{x}}}
\newcommand{\g}{{\bm{g}}}
\newcommand{\del}{\partial}
\newcommand{\f}{{\bm{f}}}
\newcommand{\aff}{{\text{aff}}}
\newcommand{\y}{{\bm{y}}}

\newcommand{\uu}{{\bm{u}}}
\newcommand{\dd}{{\bm{d}}}
\newcommand{\0}{{\bf{0}}}
\newcommand{\m}{\mathfrak{m}}

\newcommand{\p}{\mathfrak{p}}
\newcommand{\q}{\mathfrak{q}}
\newcommand{\ee}{{\bf{e}}}
\newcommand{\e}{\epsilon}

\DeclareMathOperator{\IPD}{IPD}

\DeclareMathOperator{\OCM}{\Omega CM}
\DeclareMathOperator{\Ima}{Im}
\DeclareMathOperator{\Ker}{Ker}
%\DeclareMathOperator{\com}{c}
\DeclareMathOperator{\tr}{tr}
\DeclareMathOperator{\Tr}{Tr}
\DeclareMathOperator{\Coker}{coker}
\DeclareMathOperator{\Aut}{Aut}
\DeclareMathOperator{\Inn}{Inn}
\DeclareMathOperator{\adj}{adj}
\DeclareMathOperator{\nil}{nil}
\DeclareMathOperator{\pd}{pd}
\DeclareMathOperator{\Gpd}{Gpd}



\DeclareMathOperator{\Gfd}{Gfd}
\DeclareMathOperator{\Gdim}{G-dim}
\DeclareMathOperator{\add}{add}
\DeclareMathOperator{\Add}{Add}
\DeclareMathOperator{\Rad}{Rad}
\DeclareMathOperator{\rad}{rad}
\DeclareMathOperator{\ca}{ca}
\DeclareMathOperator{\Rfd}{Rfd}
\DeclareMathOperator{\rf}{RF}
\DeclareMathOperator{\df}{def}
\DeclareMathOperator{\jac}{jac}
\DeclareMathOperator{\gldim}{gl.dim}
\DeclareMathOperator{\CM}{CM}

\DeclareMathOperator{\id}{id}
\DeclareMathOperator{\Spec}{Spec}
\DeclareMathOperator{\ac}{ac}
\DeclareMathOperator{\TF}{TF}
\DeclareMathOperator{\coh}{coh}


\DeclareMathOperator{\ann}{ann}
\newcommand{\Cat}{\mathsf{C}}
\newcommand{\Set}{\mathsf{Set}}
\DeclareMathOperator{\ob}{Obj}
\DeclareMathOperator{\Hom}{Hom}
\DeclareMathOperator{\HOM}{HOM}
\DeclareMathOperator{\Ext}{Ext}
\DeclareMathOperator{\EXT}{EXT}
\DeclareMathOperator{\End}{End}
\DeclareMathOperator{\Min}{Min}
\DeclareMathOperator{\V}{V}
\DeclareMathOperator{\Y}{Y}
\DeclareMathOperator{\het}{height}
\DeclareMathOperator{\level}{level}
\DeclareMathOperator{\maxSpec}{maxSpec}

\DeclareMathOperator{\thick}{\mathsf{thick}}
\DeclareMathOperator{\Thick}{\mathsf{Thick}}
\DeclareMathOperator{\Tor}{Tor}
\DeclareMathOperator{\Kos}{Kos}
\DeclareMathOperator{\reg}{reg}
\DeclareMathOperator{\Reg}{Reg}
\DeclareMathOperator{\Sing}{Sing}
\DeclareMathOperator{\Supp}{Supp}
\DeclareMathOperator{\Ass}{Ass}
\DeclareMathOperator{\smd}{smd}
\DeclareMathOperator{\supp}{\mathsf{supp}}
\DeclareMathOperator{\grade}{grade}
\DeclareMathOperator{\codim}{codim}

\newcommand{\shift}{{\mathsf{\Sigma}}}
\DeclareMathOperator{\RHom}{\mathsf{RHom}}

\newcommand{\ot}{\otimes^{\mathsf{L}}}
\newcommand{\xla}{\xleftarrow}
\newcommand{\xra}{\xrightarrow}
\newcommand{\DG}{\mathsf{DG}}
\newcommand{\can}{\mathsf{can}}
\newcommand{\soc}{\mathsf{soc}}
\newcommand{\CDG}{\mathsf{CDG}}
\newcommand{\per}{\mathsf{perf}}
\newcommand{\sg}{\mathsf{sg}}
\newcommand{\join}{\mathsf{Join}}
\newcommand{\Mod}{\mathsf{Mod}}
\newcommand{\Inj}{\mathsf{Inj}}
\newcommand{\mo}{\mathsf{mod}}

\newcommand{\fp}{\mathsf{fp}}

\newcommand{\Gr}{\mathsf{Gr}}
\newcommand{\gr}{\mathsf{gr}}
\newcommand{\proj}{\mathsf{proj}}
\newcommand{\inj}{\mathsf{inj}}
\newcommand{\Qcoh}{\mathsf{Qcoh}}
\newcommand{\Proj}{\mathsf{Proj}}
\newcommand{\FPD}{\mathsf{FPD}}
\newcommand{\fpd}{\mathsf{fpd}}
\newcommand{\GProj}{\mathsf{GProj}}
\newcommand{\Gproj}{\mathsf{Gproj}}
\newcommand{\GrInj}{\mathsf{GrInj}}

\newtheorem{theorem}{Theorem}[section]
\newtheorem*{Theorem}{Theorem}
\newtheorem*{Question}{Question}

\newtheorem*{Conjecture}{Conjecture}

\newtheorem*{Corollary}{Corollary}
\newtheorem{proposition}[theorem]{Proposition}

\newtheorem{fact}[theorem]{Fact}
\newtheorem{lemma}[theorem]{Lemma}
\newtheorem{corollary}[theorem]{Corollary}
\newtheorem{quest}[theorem]{Question}

\theoremstyle{definition}
\newtheorem{example}[theorem]{Example}
\newtheorem{remark}[theorem]{Remark}
\newtheorem{notation}[theorem]{Notation}
\newtheorem{definition}[theorem]{Definition}
\newtheorem*{outline}{Outline}
\newtheorem{Construction}[theorem]{Construction}
\newtheorem{chunk}[theorem]{}
\newenvironment{bfchunk}{\begin{chunk}\textbf}{\end{chunk}}
\usepackage[utf8]{inputenc}


\newtheorem*{ack}{Acknowledgements}
%\newtheorem{Thm}{Theorem}
%\newtheorem*{Prop}{Proposition}









\title[On Takahashi's questions about dominant local rings]{On Takahashi's questions about dominant local rings}
\date{\today}



\begin{document}

\begin{abstract}
This article studies dominant local rings, a notion introduced by Takahashi. We prove that, for a flat local homomorphism $R \longrightarrow S$ between commutative noetherian local rings, dominance descends from $S$ to $R$, thereby answering a question of Takahashi affirmatively. We also establish several sufficient conditions under which a Cohen--Macaulay local ring of finite CM type is dominant, supporting a conjecture of Takahashi.
\end{abstract}
\maketitle

\section{Introduction}

Let $R$ be a commutative noetherian local ring. In \cite{Takahashi:2023}, Takahashi introduced the notion of \emph{dominant local rings}. By definition, $R$ is dominant if the residue field of $R$ belongs to every nonzero thick subcategory of the singularity category of $R$; c.f. \Cref{dominant}. Dominant local rings enjoy several desirable homological properties, including being $\Tor$-friendly.  In \cite{Takahashi:2023}, Takahashi verified that several important classes of local rings are dominant. These include hypersurfaces, Cohen--Macaulay local rings with infinite residue fields and minimal multiplicity, local rings with quasi-decomposable maximal ideals introduced in \cite{NT}, and Burch rings introduced in \cite{DKT2020}. 

The classification of thick subcategories is a fundamental topic in algebra and representation theory. The classification of thick subcategories of the category of perfect complexes over a commutative noetherian ring dates back to the work of Hopkins \cite{Hopkins} and Neeman \cite{Neeman92-DR}. In \cite{Takahashi:2010}, Takahashi classified the thick subcategories of the singularity category of a hypersurface. A generalization of this result to complete intersections was established by Stevenson \cite{Stevenson}.  Recently, in \cite{Takahashi:2023}, Takahashi classified the thick subcategories of the singularity category of a local ring whose certain localizations are dominant local rings. This result highlights the importance of dominant local rings. 

In \cite{Takahashi:2026}, Takahashi introduced and studied the notion of uniformly dominant local rings. Kobayashi and Takahashi \cite{KT2026} further developed this theory. More recently, Kimura, Mifune, Otake, and Takahashi \cite{KMOT} studied the notion of quasi-dominant local rings introduced in \cite{Takahashi:2023} and showed that it is closely related to local rings over which every Gorenstein projective module is projective.

Flat local homomorphisms are among the most important morphisms in commutative algebra. Many important properties of local rings are known to descend along flat local homomorphisms. In \cite[Question 7.5]{Takahashi:2023}, Takahashi raised the following natural question:

\begin{Question}
    Let $R\longrightarrow S$ be a flat local homomorphism between commutative noetherian local rings. If $S$ is dominant, is $R$ dominant? 
\end{Question}


The first main result of this article, \Cref{T1}, shows that the above question is true. The proof makes use of the homotopy category of acyclic complexes of injective modules, which, as shown by Krause \cite{Krause2005}, is the idempotent completion of the singularity category.


\begin{theorem}\label{T1}
(See \ref{extension})  The above question of Takahashi is true. 
\end{theorem}


For a Cohen--Macaulay local ring $R$ of finite CM type, Auslander \cite{Auslander} proved that $R$ has an isolated singularity when $R$ is complete. This was extended to excellent local rings by Leuschke and Wiegand \cite{LW:2000}, and to arbitrary Cohen--Macaulay local rings by Huneke and Leuschke \cite{HL}. In \cite{EH}, Eisenbud and Herzog classified homogeneous Cohen--Macaulay rings of finite CM type. In the Gorenstein case, Buchweitz, Greuel, and Schreyer \cite{BGS}, together with Knörrer \cite{Knorrer}, classified complete Gorenstein local rings of finite CM type as simple hypersurface singularities. Recently, Cohen--Macaulay local rings of finite syzygy CM type were studied by Dey, Kimura, Liu, and Otake \cite{DKLO}.


Building on the results of \cite{Takahashi:2023}, Takahashi proposed the following conjecture in \cite[Conjecture 9.1]{Takahashi:2023}, which is a reformulation of an earlier conjecture in \cite[Conjecture 4.2]{Takahashi:2021}. Kimura, Mifune, Otake, and Takahashi \cite{KMOT} observed that a Cohen--Macaulay local ring of finite CM type admitting a canonical module is quasi-dominant.
\begin{Conjecture}
A Cohen--Macaulay local ring of finite CM type is dominant.
\end{Conjecture}



By imposing bounds on either the Hilbert--Samuel multiplicity ${\rm e}(R)$ or the codimension $\codim(R)$ of $R$, our second main result, \Cref{T2}, verifies Takahashi's conjecture for several classes of rings. In the following, the statement for the case ${\rm e}(R)=\codim(R)+1$, namely, the minimal multiplicity case, was established by Takahashi \cite{Takahashi:2023}. The proof relies on the characterizations of Cohen--Macaulay local rings of finite CM type in \cite{DKLO} and of dominant local rings in \cite{KT2026, Takahashi:2023}

\begin{theorem}\label{T2}
 (See \ref{two cases}, \ref{support}, and \ref{three}) Let $R$ be a Cohen--Macaulay local ring of finite CM type with an infinite residue field. Then $R$ is dominant provided that one of the following conditions holds.
\begin{enumerate}
    \item ${\rm e}(R)\leq 6$.
    \item $\codim(R)\leq 3$.
    \item ${\rm e}(R)\leq \codim(R)+2$.
    \item ${\rm e}(R)=\codim(R)+3$ and the type of $R$ is at least $3$.
\end{enumerate}
\end{theorem}




\section{Notation, terminology, and preliminary}
Throughout, $R$ denotes a commutative noetherian ring, and all rings considered in this article are assumed to be commutative noetherian. Let $\Mod(R)$ denote the category of all $R$-modules, and let $\mo(R)$ denote its full subcategory consisting of finitely generated $R$-modules. For each $R$-module $M$, we write $\dim(M)$ for the Krull dimension of $M$.
\begin{chunk}
\textbf{Derived category.}
Let $\D(R)$ denote the derived category of $R$-modules, which is a triangulated category with shift functor $\Sigma$. For a complex $X\in \D(R)$, the shift $\Sigma(X)$ is defined by $\Sigma(X)_i=X_{i-1}$ and $\del_{\Sigma(X)}=-\del_X$.

Let $\D^f(R)$ (resp. $\D^b(R)$) denote the full subcategory of $\D(R)$ consisting of complexes $X$ such that $\h_i(X)$ is finitely generated over $R$ for every $i\in\mathbb Z$ (resp. $\h_i(X)=0$ for $|i|\gg0$), where $\h_i(X)$ is the $i$-th homology of $X$. Set $$\D^f_b(R)\colonequals \D^f(R)\cap\D_b(R).$$  Note that a complex over $R$ is in $\D^f_b(R)$ if and only if the total homology $\bigoplus_{i\in\Z}\h_i(X)$ is finitely generated over $R$. 
\end{chunk}
\begin{chunk}
\textbf{Singularity category.}
An object $X\in\D^f_b(R)$ is called \emph{perfect} if it is isomorphic in $\D(R)$ to a bounded complex of finitely generated projective $R$-modules. We write $\per(R)$ for the full subcategory of $\D^f_b(R)$ consisting of perfect complexes. This is a triangulated subcategory of $\D^f_b(R)$. The \emph{singularity category} of $R$ is the Verdier quotient
\[
\D_{\sg}(R)\colonequals \D^f_b(R)/\per(R).
\]
This construction was introduced by Buchweitz \cite[Definition 1.2.2]{Buchweitz}; see also \cite{Orlov}.
Note that $\D_{\sg}(R)=0$ if and only if every finitely generated $R$-module has finite projective dimension. 
\end{chunk}

\begin{chunk}
\textbf{Thick subcategory.}
Let $\T$ be a triangulated category. A full subcategory $\C$ of $\T$ is called \emph{thick} if $\C$ is closed under shifts, cones, and direct summands. 
For each object $X$ in $\T$, let $\thick_\T(X)$ denote the smallest thick subcategory of $\T$ containing $X$. This can be realized as the intersection of all thick subcategories of $\T$ containing $X$. 
For example, $\per(R)$ is a thick subcategory of $\D^f_b(R)$; see \cite[Lemma 1.2.1]{Buchweitz}. Moreover, $\per(R)=\thick_{\D^f_b(R)}(R)$.

Following \cite[2.2.4]{ABIM} and \cite[Section 2]{BonVD}, the subcategory $\thick_\T(X)$ admits the following inductive description. Define
$\thick_\T^0(X)=\{0\},
$
and let $\thick_\T^1(X)$ be the smallest full subcategory of $\T$ that contains $X$ and is closed under shifts, finite direct sums, and  direct summands. For $n\geq 2$, define $\thick_\T^n(X)$ to be the full subcategory consisting of those objects $Y\in\T$ for which there exists an exact triangle
\[
Y_1\longrightarrow Y\oplus Y'\longrightarrow Y_2\longrightarrow \Sigma (Y_1)
\]
in $\T$
with
$
Y_1\in\thick_\T^1(X)$ and $
Y_2\in\thick_\T^{n-1}(X)
$.
Then
\[
\thick_\T(X)=\bigcup_{n\geq 0}\thick_\T^n(X).
\]
\end{chunk}

\begin{chunk}
    \textbf{Localizing subcategory.}
    Let $\T$ be a triangulated category admitting arbitrary coproducts. 
A full subcategory of a triangulated category $\T$ is called \emph{localizing} if it is a triangulated subcategory closed under arbitrary coproducts. Given a class of objects $S$ in $\T$, the \emph{localizing subcategory generated by $S$}, denoted by $\operatorname{Loc}_{\T}(S)$, is the smallest localizing subcategory of $\T$ containing $S$. Note that any localizing subcategory of $\T$ is thick.
\end{chunk}
 
 \begin{chunk}\label{compactly gen}
     \textbf{Compactly generated triangulated categories.}
Let $\T$ be a triangulated category admitting arbitrary coproducts. An object $X$ of $\T$ is called \emph{compact} if the functor $\Hom_{\T}(X,-)$ preserves arbitrary coproducts. Equivalently, for every family of objects $Y_i(i\in I)$ in $\T$, the canonical map
$$
\bigoplus_{i\in I}\Hom_{\T}(X,Y_i)\longrightarrow  \Hom_{\T}(X,\bigoplus_{i\in I}Y_i)
$$
is an isomorphism. Denote by $\T^{c}$ the full subcategory of $\T$ consisting of all compact objects. 

The triangulated category $\T$ is said to be \emph{compactly generated} if there exists a set of compact objects $S$ such that
$
\operatorname{Loc}_\T(S)=\T.
$
In this situation, the compact objects are precisely those lying in the thick subcategory generated by $S$; that is,
$
\T^{c}=\thick_{\T}(S),
$ 
see \cite[Lemma 2.2]{Neeman92}. For example, the derived category of $R$ is compactly generated by $R$. 
 \end{chunk}





\begin{chunk}
    \textbf{Codimension.} 
    Let $(R,\m,k)$ be a local ring. We use $\embdim(R)$ to denote the \emph{embedding dimension} of $R$, given by
$\embdim(R)\colonequals\rank_k(\m/\m^2),
$
which is the minimal number of generators of $\m$. The \emph{codimension} of $R$ is the integer
\[
\codim(R)\colonequals\rank_k(\m/\m^2)-\dim(R).
\]
\end{chunk}



\begin{chunk}
\textbf{Maximal Cohen--Macaulay module.}
For a finitely generated
$R$-module $M$, recall that $M$ is \emph{maximal Cohen--Macaulay} if
\[
\depth_{R_\p}(M_\p)\geq \dim (R_\p)
\]
for all $\p\in\Spec(R)$; here, $\depth_{R_\p}(M_\p)$ denote the depth of $M_\p$ over $R_\p$. See details in \cite[Section 2.1]{BH}.

The notation $\CM(R)$ will be used for the full subcategory of $\mo(R)$ whose
objects are the maximal Cohen--Macaulay $R$-modules. In this terminology,
$R$ is a Cohen--Macaulay ring exactly when $R\in\CM(R)$.
\end{chunk}
\begin{chunk}\label{multiplicity}
\textbf{Hilbert--Samuel multiplicity.}
Let $(R,\m)$ be a local ring. For each finitely generated $R$-module $M$, the \emph{Hilbert--Samuel multiplicity} of $M$ with respect to $\m$ is defined by
\[
{\rm e}(M)\colonequals \lim_{n\to\infty} d!\cdot\frac{\ell(M/\m^nM)}{n^d},
\]
where $d=\dim(M)$ and $\ell(-)$ represents the length of the $R$-module. See \cite[Section 4.6]{BH} for details.

Assume $k$ is infinite and $R$ is Cohen--Macaulay.  Then $\m$ contains a system of parameter $x_1,\ldots,x_t$ such that $Q\colonequals (x_1,\ldots,x_t)$ is a minimal reduction of $\m$ in the sense of Northcott and Rees \cite[Definition 2]{NR}.
Also, note that
$
{\rm e} (R)=\ell (R/Q);
$
see \cite[Exercise 4.6.24(a)]{BH}.
\end{chunk}
\begin{chunk}
\textbf{Homotopy category of (acyclic) complexes of injective modules.} Let $\K(\Inj R)$ denote the homotopy category of complexes of injective $R$-modules. Set
$$
\K_{\ac}(\Inj R)=\{X\in\K(\Inj R)\mid \h_i(X)=0 \text{ for all }i\in\Z\}.
$$
 This is a localizing subcategory of $\K(\Inj R)$.  Due to Krause \cite[Proposition 2.3]{Krause2005}, $\K(\Inj R)$ is compactly generated and there is a triangle equivalence
 $$
 \text{i}\colon \D^{f}_b(R)\xrightarrow \equiv\K(\Inj R)^{c}
 $$
 induced by taking injective resolutions.
\end{chunk}

\begin{chunk}\label{dominant}
\textbf{(Uniform) dominant local ring.}
Let $(R,\m,k)$ be a local ring. $R$ is called \emph{dominant} if
\[
k\in \thick_{\D_{\sg}(R)}(X)
\]
for every nonzero complex $X\in\D_{\sg}(R)$. See details in \cite{Takahashi:2023}. Note that the condition $k\in\thick_{\D_{\sg}(R)}(X)$ is also equivalent to that $k\in\thick_{\D^f_b(R)}(X\oplus R)$. 

Note that every object in $\D_{\sg}(R)$ is isomorphic to $\Sigma^nM$ for some finitely generated $R$-module $M$ and some $n\in\mathbb Z$. Therefore, the local ring $(R,\mathfrak m,k)$ is dominant if and only if
$k\in \thick_{\D_{\sg}(R)}(M)
$
for every nonzero finitely generated $R$-module $M$ of infinite projective dimension.

In \cite[Definition 5.3]{Takahashi:2026}, Takahashi introduced the \emph{dominant index} of $R$, defined by
\[
{\rm dx}(R)\colonequals \inf\{ n\in\mathbb{Z}_{\geq -1}\mid 
k\in\thick^{n+1}_{\D_{\sg}(R)}(X) \text{ for every }0\neq X\in\D_{\sg}(R)\}.
\]
A local ring $R$ is said to be \emph{uniformly dominant} if ${\rm dx}(R)<\infty$. Clearly, every uniformly dominant local ring is dominant.

\end{chunk}
\begin{example}\label{example}
    By \cite[Proposition 5.10]{Takahashi:2023}, known examples of dominant local rings include hypersurfaces, Cohen--Macaulay local rings with infinite residue fields and minimal multiplicity (c.f. \Cref{mini-mul}), local rings with quasi-decomposable maximal ideals \cite{NT}, and Burch rings \cite{DKT2020}.
\end{example}


\section{Descent of dominance}
The main result is \Cref{T1} from the introduction; see \Cref{extension}. It shows that the following question of Takahashi is true. 
\begin{chunk}\label{flathom}
   In \cite[Question 7.5]{Takahashi:2023},  Takahashi asked:
   \begin{Question}
       Let $(R,\m)\rightarrow (S,\n)$ be a flat local homomorphism. If $S$ is dominant, is $R$ dominant?
   \end{Question}
\end{chunk}

\begin{chunk}\label{Letz}
     Let $(R,\m)\rightarrow (S,\n)$ be a flat local homomorphism. Let $X,G$ be objects in $\D^f_b(R)$. For each $n\geq 0$, Letz \cite[Corollary 2.11]{Letz} observed that
     $$
     X\in \thick^n_{\D^f_b(R)}(G)\iff S\otimes_R X\in \thick^n_{\D^f_b(S)}(S\otimes_R G).
     $$
     In particular, $X\in\thick_{\D^f_b(R)}(G)$ if and only if $S\otimes_R X\in \thick_{\D^f_b(S)}(S\otimes_R G)$. Take $G=R$, one get
     $$
     X=0 \text{ in }\D_{\sg}(R) \iff S\otimes_R X \text{ in } \D_{\sg}(S).
     $$
\end{chunk}
For a commutative noetherian ring $A$, recall that a complex $J$ of injective $A$-modules is called \emph{$K$-injective} if the Hom complex
$
\Hom_A(E,J)
$
is acyclic for every acyclic complex $E$ of $A$-modules. 

\begin{chunk}\label{recollement}
Let $A$ be a commutative noetherian ring. By \cite[Corollary 4.3]{Krause2005}, there is a recollement
\[
\begin{tikzcd}[column sep=large, row sep=large]
\K_{\ac}(\Inj A)
  \arrow[r, "I_A"]
&
\K(\Inj A)
  \arrow[l, shift left=2.2ex, "I_{A,\rho}"]
  \arrow[l, shift right=2.2ex, "I_{A, \lambda}"']
  \arrow[r, "Q_A"]
&
\D(A)
  \arrow[l, shift left=2.2ex, "Q_{A,\rho}"]
  \arrow[l, shift right=2.2ex, "Q_{A,\lambda}"']
\end{tikzcd}
\]
in the sense of \cite{BBD}, where $I_A$ and $Q_A$ are canonical functors. Note that the right adjoint $Q_{A,\rho}$ of $Q_A$ is induced by taking $K$-injective resolutions. 
Let 
$$
\sigma_A=I_{A,\lambda} Q_{A,\rho}\colon \D(A)\rightarrow \K_{\ac}(\Inj A)
$$
 denote the \emph{stable localization} functor. This is introduced in 
 \cite[Section 5]{Krause2005}. By \cite[Corollary 5.5]{Krause2005}, the composition of the following functors
 $$
 \Mod(A)\to \D(A)\xrightarrow{\sigma_A} \K(\Inj A)
 $$
 preserves all coproducts and $\sigma_A(A)=0$. Moreover, $\sigma_A$ induces a fully faithful functor
 $$
 \overline{\sigma_A}\colon \D_{\sg}(A)\rightarrow \K(\Inj A)^c
 $$
 which is an equivalence up to direct summands; see \cite[Corollary 5.4]{Krause2005}. 

 For each $X\in\K(\Inj A)$, there exists an exact triangle
 \begin{equation}\label{decomposition}
      Q_{A,\lambda}Q_A X\to X\to I_AI_{A,\lambda}X\to \Sigma Q_{A,\lambda}Q_A X;
 \end{equation}
 see \cite[page 1135]{Krause2005}. 
 Let \(\mathcal L_A\) denote the left orthogonal of \(\K_{\ac}(\Inj A)\) in
\(\mathrm K(\operatorname{Inj}A)\). Namely, for each object $X\in\K(\Inj A)$, it is in $\mathcal L_A$ if and only if 
 $\Hom_{\K(\Inj A)}(X,Y)=0$ for all $Y\in \K_{\ac}(\Inj A)$. 
Note that $Q_{A,\lambda}Q_A X\in \mathcal L_A$ and $I_AI_{A,\lambda}X\in\K_{\ac}(\Inj R)$. Moreover, any exact triangle in $\K(\Inj A)$
$$
X_1\to X\to X_2\to \Sigma X_1
$$
with $X_1\in \mathcal L_A$ and $X_2\in \K_{\ac}(\Inj A)$ is isomorphic to the exact triangle (\ref{decomposition}).
\end{chunk}
\begin{lemma}\label{restric-fun}
Let \(R\to S\) be a flat homomorphism. The restriction of scalars induces a coproduct-preserving triangle functor
\[
U\colon\K_{\ac}(\Inj S)\longrightarrow\K_{\ac}(\Inj R).
\]
Moreover, restriction of every \(K\)-injective complex of \(S\)-modules is \(K\)-injective over \(R\).
\end{lemma}
\begin{proof}
Assume \(J\) is an injective \(S\)-module. For each \(R\)-module \(N\), there is a natural isomorphism
\[
\operatorname{Hom}_R(N,J)
\cong
\operatorname{Hom}_S(S\otimes_RN,J).
\]
Combining that $S$ is flat over $R$ and $J$ is injective over $S$, we get that $\Hom_R(-,J)$ is an exact functor, and hence \(J\) is injective as an \(R\)-module.
Since the restriction functor is exact, it takes acyclic complexes to acyclic complexes and defines the desired functor. Note that direct sums of injective modules are injective over noetherian rings. Coproducts in the relevant homotopy categories are therefore represented termwise (see \cite[Lemma 1.1]{BN}), and restriction preserves them.

Next, we prove the second statement. Let \(J\) be a \(K\)-injective complex over \(S\), and let \(E\) be an acyclic complex of \(R\)-modules. There is a natural isomorphism of Hom complexes
\[
\operatorname{Hom}_R(E,J)
\cong
\operatorname{Hom}_S(S\otimes_RE,J).
\]
Flatness makes \(S\otimes_RE\) acyclic.  Since \(J\) is \(K\)-injective over $S$, the complex on the right is acyclic. Thus, \(J\) is \(K\)-injective over \(R\).
\end{proof}

\begin{remark}
 Keep the assumption in \Cref{restric-fun}. Let $X$ be a complex of $S$-modules. We still denote it by $X$ when viewing it as a complex of $R$-modules. For $X\in \K_{\ac}(\Inj S)$, we use $U(X)$ to denote the same complex regarded as a complex of $R$-modules, which belongs to $\K_{\ac}(\Inj R)$. This notation is introduced only to make the proof of \Cref{extension} clearer and should not cause any confusion.
\end{remark}

\begin{proposition}\label{res-f}
Let \(R\to S\) be a flat homomorphism and suppose that
$\operatorname{pd}_R S<\infty. 
$
Then

(1) For each $X\in\D(S)$,  
there is an isomorphism
$
U\sigma_S(X)\cong\sigma_R(X)
$
in $\K(\Inj R)$. 

(2)
If
$
(R,\mathfrak m,k)\longrightarrow(S,\mathfrak n,l)
$
is local, then 
$
\sigma_R(k)
$
is a direct summand of $U\sigma_S(l)$ 
in \(\K_{\ac}(\Inj R)\).
\end{proposition}
\begin{proof}
(1) Let \(E\) be an acyclic complex of injective \(R\)-modules. Since \[
\operatorname{Hom}_S(-,\operatorname{Hom}_R(S,E^j))
\cong
\operatorname{Hom}_R(U(-),E^j),
\]
we get that the complex
$
\operatorname{Hom}_R(S,E)
$
is a complex of injective modules over \(S\).
Fix $i\in\Z$, let $K_i$ denote the kernel of the map $\del_i^E\colon E_i\to E_{i-1}$. Note that for each $n\geq 0$,  $K_i$ is the $n$-th cosyzygy of  $\Omega_{-{n}}^R(K_{i+n})$. Combining with this, we get
$$
\Ext^1_R(S,K_i)\cong \Ext^{1+n}_R(S,K_{i+n})
$$
for each $n\geq 0$.  This yields that $\Ext^1_R(S,K_i)=0$ as $\pd_R(S)<\infty$. 
Since this holds for each $i\geq 0$, \(\operatorname{Hom}_R(S,E)\) is an acyclic complex of injective \(S\)-modules.

 Next, keep the notation as in \Cref{recollement}. If \(L\in\mathcal L_S\) and
\(E\in\K_{\ac}(\Inj R)\), there is an adjunction isomoprhism
\[
\operatorname{Hom}_{\mathrm K(\operatorname{Inj}R)}(L,E)
\cong
\operatorname{Hom}_{\mathrm K(\operatorname{Inj}S)}
\bigl(L,\operatorname{Hom}_R(S,E)\bigr).
\]
This is zero as we have $\Hom_R(S,E)\in \K_{\ac}(\Inj S)$. 
Thus, the restriction of the object in $
\mathcal L_S$ is in  $\mathcal L_R.$

For each \(X\in\mathrm D(S)\), then $Q_{S,\rho}(X)\in \K(\Inj S)$,  and the recollement in \Cref{recollement} induces an exact triangle
\[
L_X\longrightarrow Q_{S,\rho}(X)
\longrightarrow I_S I_{S,\lambda} Q_{S,\rho}(X)
\longrightarrow\Sigma L_X
\]
in $\K(\Inj S)$,
where \(L_X\in\mathcal L_S\); see \Cref{recollement}. Note that $I_S I_{S,\lambda} Q_{S,\rho}(X)=\sigma_S(X)$ is acyclic. After restriction, the first term is in \(\mathcal L_R\), the middle term is in $\K(\Inj R)$ by the proof of \Cref{restric-fun}, and the third one is in $\K_{\ac}(\Inj R)$ by \Cref{restric-fun}. The uniqueness of the corresponding orthogonal decomposition in \Cref{recollement} implies
\[
U\sigma_S(X)\cong I_{R,\lambda}Q_{S,\rho}X.
\]
Since $Q_{S,\rho}X$ is $K$-injective over $S$ (see \Cref{recollement}), it follows from \Cref{restric-fun} that it is also $K$-injective over $R$ . Thus, $Q_{S,\rho}X\cong Q_{R,\rho}X$  in $\K(\Inj R)$, and hence we conclude that $U\sigma_S(X)\cong \sigma_R(X)$.

(2) The $R$-module $l$ is also a $k$-vector space. Thus, $k$ is a direct summand of $l$ over $k$.  This also yields that $k$ is a direct summand of $l$ as $R$-modules. Applying the functor \(\sigma_R\), we get that the module $\sigma_R(k)$ is a direct summand of $\sigma_R(l)$. 
Combining with $\sigma_R(l)\cong U\sigma_S(l),
$
we conclude that $\sigma_R(k)$ is a direct summand of $U\sigma_S(l)$ in $\K_{\ac}(\Inj R)$. 
\end{proof}

\begin{proposition}\label{res:contain}
Let \(R\to S\) be a flat homomorphism with \(\operatorname{pd}_RS<\infty\), and let \(M\) be a finitely generated \(R\)-module. Then
$
\sigma_R(S\otimes_RM)
\in
\operatorname{Loc}_{\K_{\ac}(\Inj R)}(\sigma_R(M)) 
$ and we have
\[
U\sigma_S(S\otimes_RM)
\in
\operatorname{Loc}_{\K_{\ac}(\Inj R)}(\sigma_R(M)).
\]
\end{proposition}
\begin{proof}
Put \(d=\operatorname{pd}_RS\), and choose a projective resolution
\[
0\longrightarrow P_d\longrightarrow\cdots
\longrightarrow P_0\longrightarrow S\longrightarrow0.
\]
Because \(S\) is flat, every syzygy in this resolution is flat.  Combining with this, by applying the functor \(-\otimes_RM\), we get  a long exact sequence
\begin{equation}\label{equation}
    0\longrightarrow P_d\otimes_RM\longrightarrow\cdots
\longrightarrow P_0\otimes_RM
\longrightarrow S\otimes_RM\longrightarrow0.
\end{equation}

For every set \(I\), it follows from \Cref{recollement} that
$
\sigma_R(M^{(I)})\cong\sigma_R(M)^{(I)}.
$
%Indeed, choose a bounded-below injective resolution \(J_M\) of \(M\). Since \(R\) is noetherian, \(J_M^{(I)}\) is again a bounded-below complex of injectives and is a \(K\)-injective resolution of \(M^{(I)}\). The left adjoint used in stabilization preserves coproducts, proving the displayed isomorphism.
Since each \(P_i\) is a direct summand of a free module \(R^{(I_i)}\), the module \(P_i\otimes_RM\) is a direct summand of \(M^{(I_i)}\). Combining with $
\sigma_R(M^{(I)})\cong\sigma_R(M)^{(I)}
$, we get
\begin{equation}\label{123}
   \sigma_R(P_i\otimes_RM)
\in
\operatorname{Loc}_{\K_{\ac}(\Inj R)}(\sigma_R(M))
\end{equation}
for each $0\leq i\leq d$. 
By (\ref{equation}), $S\otimes_R M\in \thick_{\D(R)}(\oplus_{i=0}^d P_i\otimes _R M)$, and hence 
\begin{equation}\label{234}
    \sigma_R(S\otimes_R M)
\in \thick_{\K_{\ac}(\Inj R)}(\oplus_{i=0}^d \sigma_R(P_i\otimes _R M));
\end{equation}
see \cite[Lemma 2.4(6)]{ABIM}.
By (\ref{123}) and (\ref{234}), we have 
\[
\sigma_R(S\otimes_RM)
\in
\operatorname{Loc}_{\K_{\ac}(\Inj R)}(\sigma_R(M)).
\]
The second statement follows by combining this with \Cref{res-f}.
\end{proof}

\begin{theorem}\label{extension}
Let
$
(R,\mathfrak m,k)\longrightarrow(S,\mathfrak n,l)
$
be a flat local homomorphism.
If \(S\) is dominant, then \(R\) is dominant. 
\end{theorem}
\begin{proof}
By \cite[Theorem 3.2.6]{RG} and \cite[Proposition 6]{Jensen}, every flat module has finite projective dimension over a commutative noetherian ring of finite Krull dimension; see also \cite[Proposition 5.9]{DKLO}. In particular, $\pd_R(S)<\infty$. Hence, the assumptions in \Cref{res-f} and \Cref{res:contain} are satisfied.

Fix a finitely generated \(R\)-module \(M\) of infinite projective dimension. By \Cref{Letz}, we have
$
\operatorname{pd}_S(S\otimes_RM)=\infty.
$
Together with the dominance of $S$, this yields
\[
l\in
\operatorname{thick}_{\D^f_b(S)}
(S\oplus (S\otimes_RM)).
\]
Since $\sigma_S(S)=0$ (see \Cref{recollement}), we obtain
$
\sigma_S(l)\in
\operatorname{thick}_{\K_{\ac}(\Inj S)}
\bigl(\sigma_S(S\otimes_RM)\bigr).
$
After restriction, \Cref{res:contain} implies that
$
U\sigma_S(l)\in
\operatorname{Loc}_{\K_{\ac}(\Inj R)}(\sigma_R(M)).
$
By \Cref{res-f}, the compact object \(\sigma_R(k)\) is a direct summand of \(U\sigma_S(l)\). Therefore,
$$
\sigma_R(k) \in
\operatorname{Loc}_{\K_{\ac}(\Inj R)}(\sigma_R(M)).
$$

Note that category \(\K_{\ac}(\Inj R)\) is compactly generated (see \Cref{compactly gen}), and the stabilizations of finitely generated modules are compact (see \Cref{recollement}). Hence, we obtain the first inclusion below:
\begin{align*}
   \sigma_R(k) &\in
\K_{\ac}(\Inj R)^c\cap \operatorname{Loc}_{\K_{\ac}(\Inj R)}
 (\sigma_R(M))\\
 & \subseteq \operatorname{Loc}_{\K_{\ac}(\Inj R)}
 (\sigma_R(M))^c\\
& =\thick_{\K_{\ac}(\Inj R)^c}(\sigma_R(M)), 
\end{align*}
where the equality follows by combining \Cref{compactly gen} with the fact that $\K(\Inj R)^c$ is a thick subcategory of $\K_{\ac}(\Inj R)$.

Combining this with the fully faithful functor $\D_{\sg}(R)\to \K_{\ac}(\Inj R)$ induced by $\sigma_R$, we conclude from \cite[Proposition 3.15]{DLL} that
\[
k\in
\operatorname{thick}_{\D_{\sg}(R)}(M).
\]
This proves that $R$ is dominant, since $M$ was chosen as an arbitrary finitely generated $R$-module of infinite projective dimension.
\end{proof}

\begin{remark}\label{completion}
       For a local ring $(R,\m)$, Takahashi \cite[Corollary 5.8]{Takahashi:2023} proved that $R$ is dominant if and only if $\widehat{R}$ is dominant. Note the descent direction is also a direct consequence of \Cref{extension}. 
\end{remark}

\section{Cohen--Macaulay ring with finite CM type}
Throughout this section, $(R,\m,k)$ will be a Cohen--Macaulay local ring.


   








\begin{chunk}
Recall that a local ring is said to have \emph{finite CM type} if it has only finitely many isomorphism classes of indecomposable maximal Cohen--Macaulay modules. In \cite[Conjecture 9.1]{Takahashi:2023}, Takahashi proposed the following conjecture.
\begin{Conjecture}
A Cohen--Macaulay local ring of finite CM type is dominant.
\end{Conjecture}
\end{chunk}


\begin{chunk}
Let $
\mathbf{P}=\cdots \longrightarrow P_{1}\stackrel{\partial_{1}}\longrightarrow
P_0\stackrel{\partial_0}\longrightarrow P_{-1}\longrightarrow\cdots
$
be a complex of projective $R$-modules. The complex $\mathbf P$ 
is called \emph{totally acyclic} if it is acyclic and
$\Hom_R(\mathbf{P},F)$ is acyclic for any projective $R$-module $F$.

An $R$-module $M$ is called \emph{Gorenstein projective} provided that there exists
a totally acyclic complex $\mathbf{P}$ such that
$
M\cong\operatorname{Im}(\partial_0).
$
See details in \cite[Chapter 10]{EJ}.
Let $\GProj(R)$ denote the category of all Gorenstein projective $R$-modules, and let $\Proj(R)$ denote the category of all projective $R$-modules. Then $\Proj(R)\subseteq \GProj(R)$.  

The ring $R$ is said to be \emph{$\G$-regular} if $\GProj(R)\cap \mo(R)=\Proj(R)\cap \mo(R)$. See details in \cite{Takahashi2008}.
\end{chunk}

\begin{chunk}\label{Tor-friendly}
Let $R$ be a commutative noetherian ring.    
$R$ is said to be \emph{$\Tor$-friendly} if, for every $M,N\in\mo(R)$ satisfying
$\Tor^R_i(M,N)=0$ for $i\gg 0$, either $M$ or $N$ has finite projective dimension.
$R$ is said to be \emph{$\Ext$-friendly} if, for every $M,N\in\mo(R)$ satisfying
$\Ext^i_R(M,N)=0$ for $i\gg0$, either $M$ has finite projective dimension or $N$ has finite injective dimension.

In \cite[Proposition 5.5]{AINSW}, Avramov, Iyengar, Nasseh, and Sather-Wagstaff established that every $\Tor$-friendly local ring is $\Ext$-friendly. By \cite[Theorem 6.7]{Takahashi:2023}, every dominant local ring is $\Tor$-friendly, and thus $\Ext$-friendly.
\end{chunk}
 

\begin{proposition}\label{two cases}
    Let $(R,\m,k)$ be a Cohen--Macaulay local ring of finite CM type.  Then $R$ is dominant if one of the following conditions holds.
    \begin{enumerate}
        \item  $k$ is infinite and ${\rm e}(R)\leq 6$.

        \item  $\codim(R)\leq 3$. 
    \end{enumerate}
\end{proposition}



\begin{proof}
    Since $R$ has finite CM type,  it follows from \cite[Theorem 1.3]{DKLO} that either $R$ is a hypersurface or $\GProj(R)=\Proj(R)$.
    Suppose first that $R$ is a hypersurface. Then $R$ is dominant by \Cref{example}.

    Next, assume that $\GProj(R)=\Proj(R)$. Then $R$ is $\mathsf G$-regular. If $R$ is Gorenstein, then $\Omega^d_R(k)$ is Gorenstein projective, and hence projective, where $d=\dim(R)$. It follows that $R$ is regular, and thus dominant. If $R$ is not Gorenstein,  next we prove that $R$ is dominant under the assumption of (1) or (2).

    
 (1) Assume $R$ satisfies the assumption of (1), then \cite[Proposition 6.7(4)]{KT2026} yields that $R$ is dominant, where we use that $R$ is $\mathsf G$-regular, $k$ is infinite, and ${\rm e}(R)\leq 6$. 
 
 (2) Assume $R$ satisfies the assumption of (2). By \cite[Theorem 1.2]{DKLO}, $R$ has finite CM type if and only if the Cohen--Macaulay ring $\widehat{R}$ has finite CM type. 
 Combining this with \Cref{completion}  and $\codim(R)=\codim(\widehat{R})$, we may assume $R$ is complete. We continue under this assumption. 
Since $R$ is $\mathsf G$-regular, $R$ has no embedded deformation; see \cite[Remark 5.5(5)]{KT2026}. By \cite[Theorem 3.8]{LM2020}, $R$ is $\Ext$-friendly.  Combining this with \cite[Proposition 6.6(1)]{Takahashi:2023}, $R$ is dominant.  
\end{proof}
\begin{chunk}\label{mini-mul}
In \cite{Abhyankar}, Abhyankar proved that the Hilbert--Samuel  multiplicity of $(R,\m)$
satisfies the inequality
\[
{\rm e}(R)\geq \embdim(R)-\dim(R)+1.
\]
A local ring $R$ is said to have \emph{minimal multiplicity} if equality holds
in the above inequality.
If the residue field $k$ is infinite, then $R$ has minimal multiplicity if and
only if
\[
\m^2=Q\m
\]
for some minimal reduction $Q$ of $\m$; see 
\cite[Exercise 4.6.14]{BH}.
\end{chunk} 


\begin{theorem}\label{support}
    Let $(R,\m,k)$ be a Cohen--Macaulay local ring of finite CM type. Assume that $k$ is infinite and that ${\rm e}(R)=\embdim(R)-\dim(R)+2$. Then $R$ is dominant. 
\end{theorem}
\begin{proof}
The same proof in \Cref{two cases} shows that $R$ is either a hypersurface or $\GProj(R)=\Proj(R)$. This yields that $R$ is dominant if $R$ is a hypersurface or Gorenstein with $\GProj(R)=\Proj(R)$. Next, we assume $R$ is not Gorenstein with $\GProj(R)=\Proj(R)$. 

   Let $d=\dim(R)$. Since $k$ is infinite, $\m$ contains a system of parameters $x_1,\ldots,x_d$ of $R$ such that $Q=(x_1,\ldots,x_d)$ is a minimal reduction of $\m$. Then \Cref{multiplicity} yields  
   $$
   {\rm e}(R)=\ell(R/Q)=\embdim(R/Q)+2. 
   $$
   Set $A\colonequals R/Q$. Then $\n\colonequals \m/Q$ is the maximal ideal of $A$. By the above equality, we have $\sum_{i\geq 2}\rank_k(\n^i/\n^{i+1})=1$. This yields that $\rank_k(\n^2/\n^3)=1$. Combining with that $A$ is non-Gorenstein artinian, it follows from \cite[Proposition 6.3(1)(c)]{KT2026} that ${\rm dx}(A)\leq 0$; see \Cref{dominant} for the definition of the dominant index ${\rm dx}(A)$. Thus, \cite[Lemma 6.1(3)]{KT2026} yields that
   $$
   {\rm dx}(R)\leq (d+1)({\rm dx}(A)+1)-1=d.
   $$
   Hence, $R$ is uniformly dominant, and hence dominant. 
\end{proof}
\begin{remark}
A Cohen--Macaulay local ring $(R,\m,k)$ with ${\rm e}(R)=\embdim(R)-\dim(R)+1$ precisely means that $R$ has minimal multiplicity. Such rings are always dominant when $k$ is infinite; see \Cref{example}.
\end{remark}
Let ${\rm r}(R)$ denote the \emph{type} of the local ring $(R,\m,k)$. By definition,
\[
{\rm r}(R)\colonequals \rank_k\Ext^t_R(k,R),
\]
where $t=\depth(R)$.


\begin{chunk}\label{test}
   Let $(R,\m)$ be a local ring, and assume that $x\in\m$ is an $R$-regular element. Takahashi \cite[Theorem 5.6]{Takahashi:2023} observed that if $R/(x)$ is dominant, then so is $R$. Moreover, it was shown in loc. cit. that the converse holds provided that $x\notin\m^2$.
\end{chunk}
\begin{proposition}\label{three}
      Let $(R,\m,k)$ be a Cohen--Macaulay local ring of finite CM type. Assume that $k$ is infinite, ${\rm e}(R)=\embdim(R)-\dim(R)+3$, and ${\rm r}(R)\geq 3$. Then $R$ is dominant. 
\end{proposition}
\begin{proof}
    By the proof of \Cref{support}, it remains to prove the desired statement when $R$ is not Gorenstein and $\GProj(R)=\Proj(R)$.   Let $d=\dim(R)$. The same reason as in \Cref{support}, $\m$ contains a system of parameters $x_1,\ldots,x_d$ of $R$ such that $Q=(x_1,\ldots,x_d)$ is a minimal reduction of $\m$. Then \Cref{multiplicity} yields that
   $$
   {\rm e}(R)=\ell(R/Q)=\embdim(R/Q)+3. 
   $$
   Set $A\colonequals R/Q$. Then $\n\colonequals \m/Q$ is the maximal ideal of $A$. It follows from  the above equality that $\sum_{i\geq 2}\rank_k(\n^i/\n^{i+1})=2$. Thus, there are two possible cases for $A$:  
   \begin{enumerate}
       \item $\rank_k(\n^2/\n^{3})=1$ and $\n^3\neq 0$.

       \item $\rank_k(\n^2/\n^3)=2$ and $\n^3=0$. 
   \end{enumerate}
   
If $A$ is of the form (1), then $\rank_k(\n^2/\n^{3})=1$. The same argument as in the proof of \Cref{support} yields that $R$ is dominant.

If $A$ is of the form (2), then $\n^2\subseteq \soc(A)$, where $\soc(A)$ denotes the socle of $A$. By \cite[Lemma 3.1.16]{BH}, ${\rm r}(A)={\rm r}(R)\geq 3$.  Combining with ${\rm r}(A)=\rank_k(\soc(A))$, we get $\n^2\neq \soc(A)$. Choose $z\in \soc(A)\setminus \n^2$. Next, we claim that $\n=Az\oplus J$ as $A$-modules for some ideal $J$ of $A$. Indeed, consider the natural embedding $Az\to \n$. This is a split injection. In fact, since $Az$ is a simple module and it is a direct summand of $\n/\n^2$, there exists a split surjection $g\colon \n/\n^2\to Az$ sending $\overline{z}$ to $z$. Hence, the composition
$
\n\rightarrow \n/\n^2\xrightarrow{g}Az
$
is a retraction of the natural embedding $Az\to \n$. Thus, $\n$ is a decomposable maximal ideal. By \Cref{example}, $A$ is dominant. Since $x_1,\ldots,x_d$ is an $R$-regular sequence, \Cref{test} yields that $R$ is also dominant. 
\end{proof}
\begin{ack}
  This work was carried out by the author with the assistance of Eureka, a multi-agent system developed by JIUCHONG at the University of Science and Technology of China for mathematical research. The author would like to thank the JIUCHONG team for providing access to Eureka. The author was supported by the National Natural Science Foundation of China (No. 12401046).
\end{ack}
\bibliographystyle{amsplain}
\bibliography{ref}
\end{document} 


       